\par\noindent\textbf{(1)} 微小电流环的磁矩 \(\boldsymbol m=IA\hat{\boldsymbol n}\)。对各线段磁力 \(d\boldsymbol F=I\,d\boldsymbol l\times\boldsymbol B\) 求力矩，均匀场的净力为零而总力矩 \(\boldsymbol\tau=\boldsymbol m\times\boldsymbol B\)。因此取向势能 \(U=-\boldsymbol m\cdot\boldsymbol B=-m_0B\cos\theta\)。
\par\noindent\textbf{(2)} 对沿场轴非均匀强度 \(B(z)\) 且磁矩方向暂定，轴向力 \(F_z=-\partial_z U=m_0\cos\theta\,B'(z)\)。一般矢量近似为 \(\boldsymbol F=\nabla(\boldsymbol m\cdot\boldsymbol B)\)，其中磁矩取固定而非与位置任意独立变化。最大吸引力来自磁矩顺场并沿场增强方向。
\par\noindent\textbf{(3)} 若偶极环有关于转向轴的转动惯量 \(I_r\)，在固定位置 \(B>0\) 时角方程 \(I_r\ddot\theta=-m_0B\sin\theta\)。顺场 \(\theta=0\) 附近有 \(\ddot\theta+(m_0B/I_r)\theta=0\)，小振动频率 \(\sqrt{m_0B/I_r}\)；反向 \(\theta=\pi\) 的线性化则为指数不稳定。
\par\noindent\textbf{(4)} 若 \(B'(z)=0\) 但 \(B\ne0\)，轴向净力为零而取向力矩仍可非零。反转场方向可交换顺场与逆场平衡；若外场强度在空间变化太快，有限尺寸电流环必须引入高阶磁多极项。
